\section{Time stepping: the discrete energy identity}
\label{sec:energy}

\noindent\emph{Status: draft -- the update, the invariant and the stability bound are implemented and covered by the
tests whose numbers are quoted below; the dispersion relation is derived here but is not yet the subject of a
test of its own.}

The Maxwell grid equations \eqref{eq:faraday} and \eqref{eq:ampere} of Section~\ref{sec:theory} are two coupled
first-order equations in time. Integrating them is therefore not a question of solving a linear system but of
choosing an integrator that respects the structure the spatial discretisation was built to preserve. This section
states the scheme the code actually runs, derives the quadratic form it preserves exactly -- which is not the
obvious one -- and reports the measured drift. It then treats the two limits that constrain every practical run:
the stability bound on the time step, which is fixed by the smallest cell anywhere in the grid, and the dispersion
of a wavefront on the grid, which makes it travel a little slower than light. The identity of
Section~\ref{sec:invariant} is the discrete counterpart of the conservation law proved for the finite integration
technique by \cite{schuhmann2001}; the stability and dispersion results are the classical ones for this staggering
\cite{taflove2005,yee1966}.

\subsection{The update as implemented}
\label{sec:update}

The scheme is the leapfrog (Verlet) pair: the magnetic field is advanced over a half step using the electric field
at integer times, the electric field is then advanced over a full step using the magnetic field at the half step,
and the conductor constraint is re-imposed. In the notation of Section~\ref{sec:theory},
\begin{align}
  \Mm\,\frac{\mathbf{h}^{n+1/2}-\mathbf{h}^{n-1/2}}{\Delta t} &= -\Cmat\,\mathbf{e}^{n},
  \label{eq:leapfrog-h}\\
  \Em\,\frac{\mathbf{e}^{n+1}-\mathbf{e}^{n}}{\Delta t} &= \tilde\Cmat\,\mathbf{h}^{n+1/2}
  - \mathbf{j}^{n+1/2},
  \label{eq:leapfrog-e}
\end{align}
with $\tilde\Cmat=\Cmat^{\!\top}$, followed by the projection
\begin{equation}
  \hat e_i^{\,n+1} = 0 \qquad \text{for every edge } i \text{ tangential to a perfect conductor},
  \label{eq:leapfrog-pec}
\end{equation}
which is \eqref{eq:pec} applied to the solution rather than to the solution space. The stars in
\eqref{eq:leapfrog-h} and \eqref{eq:leapfrog-e} are the diagonal Hodge matrices of Section~\ref{sec:hodge}, with
$[\Em]_{ii}=\eps_{r,i}\eps_0A_{\tilde i}/L_i$ in farads and $[\Mm]_{ff}=\mu_{r,f}\mu_0A_f/L_{\tilde f}$ in henries,
as in \eqref{eq:stars}. Two of their properties do all the work below. Both are diagonal and positive for positive
material parameters, so a half step costs two sparse products by $\Cmat$ and $\tilde\Cmat$ and four elementwise
scalings, and never a linear solve. And both scale \emph{linearly} with the cell size: the ratio $A_{\tilde i}/L_i$
has dimension of length, so doubling every cell doubles every entry of $\Em$ and of $\Mm$ exactly, a fact the test
suite asserts for that factor. The stars are Hodge stars, not mass matrices in the length-invariant sense usually
meant by that name, and it is precisely this linear scaling that later produces the linear-in-$h$ bound of
Section~\ref{sec:cfl}.

The projection \eqref{eq:leapfrog-pec} is not a repair of a solution that has gone wrong. By \eqref{eq:pec} it is
a constraint on the unknown space, and the rows of $\Cmat$ belonging to pinned edges could have been deleted in the
first place; we keep them for bookkeeping and zero the result, which is equivalent because a zero column of
$\Cmat$ contributes nothing to $\Cmat\mathbf{e}$. The argument of Section~\ref{sec:invariant} therefore runs
unchanged on the restricted space, bit for bit, as promised in Section~\ref{sec:theory}.

\subsection{The energy that the scheme does not conserve}

The obvious candidate for a conserved quantity is the field energy
\begin{equation}
  U_{\mathrm{naive}} = \tfrac12\,\mathbf{e}\cdot(\Em\mathbf{e}) + \tfrac12\,\mathbf{h}\cdot(\Mm\mathbf{h})
  \ \ [\mathrm{J}],
  \label{eq:naive}
\end{equation}
which is \eqref{eq:energies}. It is \emph{not} conserved by \eqref{eq:leapfrog-h}--\eqref{eq:leapfrog-e}. To see
why, take the inner product of \eqref{eq:leapfrog-h} with $\tfrac{\Delta t}{2}(\mathbf{h}^{n+1/2}+
\mathbf{h}^{n-1/2})$ and of \eqref{eq:leapfrog-e} with $\tfrac{\Delta t}{2}(\mathbf{e}^{n+1}+\mathbf{e}^{n})$, and
use the symmetry of $\Em$ and $\Mm$. The two identities
\begin{align}
  \tfrac12\,\mathbf{e}^{n+1}\!\cdot\Em\mathbf{e}^{n+1} - \tfrac12\,\mathbf{e}^{n}\!\cdot\Em\mathbf{e}^{n}
  &= \tfrac{\Delta t}{2}\,\Cmat(\mathbf{e}^{n+1}+\mathbf{e}^{n})\cdot\mathbf{h}^{n+1/2},
  \label{eq:telescope-e}\\
  \tfrac12\,\mathbf{h}^{n+1/2}\!\cdot\Mm\mathbf{h}^{n+1/2}
  - \tfrac12\,\mathbf{h}^{n-1/2}\!\cdot\Mm\mathbf{h}^{n-1/2}
  &= -\tfrac{\Delta t}{2}\bigl(\Cmat\mathbf{e}^{n}\bigr)\cdot(\mathbf{h}^{n+1/2}+\mathbf{h}^{n-1/2})
  \label{eq:telescope-h}
\end{align}
follow, and adding them cancels the shared term $\tfrac{\Delta t}{2}(\Cmat\mathbf{e}^{n})\cdot\mathbf{h}^{n+1/2}$
only once. What survives, writing $U_{\mathrm{naive}}^{n}$ for \eqref{eq:naive} evaluated with $\mathbf{e}^{n}$ and
$\mathbf{h}^{n-1/2}$, is
\begin{equation}
  U_{\mathrm{naive}}^{n+1} - U_{\mathrm{naive}}^{n}
  = \tfrac{\Delta t}{2}\Bigl[\bigl(\Cmat\mathbf{e}^{n+1}\bigr)\cdot\mathbf{h}^{n+1/2}
  - \bigl(\Cmat\mathbf{e}^{n}\bigr)\cdot\mathbf{h}^{n-1/2}\Bigr],
  \label{eq:naive-drift}
\end{equation}
which does not vanish: the right-hand side is the one-step change of the cross term
$\tfrac{\Delta t}{2}(\Cmat\mathbf{e})\cdot\mathbf{h}$. The defect is first order in $\Delta t$, of relative size
$\omega\Delta t$. That is negligible while the step resolves the period and of order unity for a mode at the
stability limit, and Section~\ref{sec:evidence} measures the second case at $3.71\times10^{-1}$. The naive energy
is not conserved to a small error; it is the wrong quantity.

\subsection{The exact invariant}
\label{sec:invariant}

\begin{proposition}[Exact discrete energy identity]
\label{prop:invariant}
Let $\mathbf{j}=\mathbf{0}$. Then \eqref{eq:leapfrog-h} and \eqref{eq:leapfrog-e} preserve
\begin{equation}
  U^{n} = \tfrac12\,\mathbf{e}^{n}\cdot\Em\mathbf{e}^{n}
        + \tfrac12\,\mathbf{h}^{n-1/2}\cdot\Mm\mathbf{h}^{n-1/2}
        - \tfrac{\Delta t}{2}\bigl(\Cmat\mathbf{e}^{n}\bigr)\cdot\mathbf{h}^{n-1/2}
  \ \ [\mathrm{J}]
  \label{eq:invariant}
\end{equation}
identically, $U^{n+1}=U^{n}$, and the same form with every magnetic index shifted by half a step.
\end{proposition}

\begin{proof}
Add \eqref{eq:telescope-e} and \eqref{eq:telescope-h}. The right-hand side is
$\tfrac{\Delta t}{2}[(\Cmat\mathbf{e}^{n+1})\cdot\mathbf{h}^{n+1/2}-(\Cmat\mathbf{e}^{n})\cdot
\mathbf{h}^{n-1/2}]$, which by \eqref{eq:naive-drift} is exactly the increment that the third term of
\eqref{eq:invariant} subtracts. Hence $U^{n+1}-U^{n}=0$.
\end{proof}

The cross term is not an embellishment added to make the numbers look better; it is what remains of the energy
when the two fields are sampled half a step apart, and \eqref{eq:naive-drift} is the statement that omitting it
costs a systematic first-order offset. Two readings make it natural. The first is Hamiltonian:
\eqref{eq:leapfrog-h} advances $\mathbf{h}$ while $\mathbf{e}$ is frozen, \eqref{eq:leapfrog-e} advances
$\mathbf{e}$ while $\mathbf{h}$ is frozen, and each of those shears of phase space is a symplectic map; their
composition is symplectic, and \eqref{eq:invariant} is the quadratic form that the composed linear map preserves.
The second is variational: \eqref{eq:invariant} is the Hamiltonian of the split system evaluated at the staggered
midpoint $\tfrac12(\mathbf{e}^{n}+\mathbf{e}^{n+1})$, $\tfrac12(\mathbf{h}^{n-1/2}+\mathbf{h}^{n+1/2})$, which is
why the two fields appear at times differing by $\Delta t/2$.

Symplecticity is also the reason the drift is a rounding floor rather than a growing trend. A symplectic linear
map preserves \eqref{eq:invariant} up to the rounding committed in evaluating it -- two sparse products and four
scalings per step -- so each step contributes a defect of order $\epsilon_{\mathrm{mach}}U$ with no secular
amplification, in contrast with a truncation term of order $\Delta t^{k}U$, which would accumulate linearly in the
step count. Over a few thousand steps the two behaviours are separated by many orders of magnitude, which is the
whole point of Section~\ref{sec:evidence}.

With a source the identity acquires a right-hand side, and the right-hand side is the one the physics asks for:
\begin{equation}
  U^{n+1}-U^{n} = -\Delta t\,\mathbf{j}^{n+1/2}\cdot\bar{\mathbf{e}},
  \qquad \bar{\mathbf{e}} = \tfrac12\bigl(\mathbf{e}^{n}+\mathbf{e}^{n+1}\bigr),
  \label{eq:source-balance}
\end{equation}
so the injected power is $\mathbf{e}\cdot\mathbf{j}$ at the midpoint, which is \eqref{eq:power}. The energy
balance of the discrete system is thus the discrete power balance of the continuous one, with no extra term. The
only place the correspondence needs care is the hard source of Section~\ref{sec:theory}, which overrides
\eqref{eq:leapfrog-e} on one edge and therefore performs work that \eqref{eq:source-balance} does not see.

\subsection{Numerical evidence}
\label{sec:evidence}

The experiment is a $12\times10\times8$ grid, a random initial field, no source, and 500 steps. The quantity
tabulated is the relative drift
\begin{equation}
  \delta U = \max_{0\le n\le N}\frac{|U^{n}-U^{0}|}{|U^{0}|}
  \label{eq:drift}
\end{equation}
of the form under test. Table~\ref{tab:candidates} reports four candidates: the form \eqref{eq:invariant}, and the
three nearest variants obtained from it by flipping the sign of the cross term, by pairing the magnetic field of
the cross term with $\mathbf{h}^{n+1/2}$ while the energy uses $\mathbf{h}^{n-1/2}$, and by dropping the term
altogether, which recovers \eqref{eq:naive}.

\begin{table}[H]
  \centering
  \caption{Relative drift \eqref{eq:drift} of four candidate invariants over 500 undriven steps on a
  $12\times10\times8$ grid with a random initial field. The candidates differ in one term each.}
  \label{tab:candidates}
  \begin{tabular}{@{}llr@{}}
    \toprule
    Candidate & Distinguishing term & $\delta U$ \\
    \midrule
    \eqref{eq:invariant}   & $-\tfrac{\Delta t}{2}(\Cmat\mathbf{e})\cdot\mathbf{h}$ & $1.298\times10^{-15}$ \\
    wrong sign             & $+\tfrac{\Delta t}{2}(\Cmat\mathbf{e})\cdot\mathbf{h}$ & $4.79\times10^{-1}$ \\
    mixed half-step        & $\mathbf{h}^{n+1/2}$ against $\mathbf{h}^{n-1/2}$      & $3.40\times10^{-1}$ \\
    no cross term, \eqref{eq:naive} & none                                          & $3.71\times10^{-1}$ \\
    \bottomrule
  \end{tabular}
\end{table}

The separation is fourteen orders of magnitude, and it is the same separation for each of the three variants: a
sign error, a half-step mismatch, or the omission of the term each cost the full factor. The measurement therefore
does not merely confirm the form we chose, it discriminates it against its three nearest neighbours, which is a
much stronger statement than a small residual on its own. Longer runs behave identically. Over 1500 undriven steps
the demonstration script reports $8.45\times10^{-16}$, and a live run of the conservation test reports
$1.669\times10^{-15}$; the three figures differ by the rounding of different data over different step counts, not
by a trend, which is what a rounding floor looks like when it is sampled three times. For scale,
$10^{-15}$ relative is only a few multiples of $\epsilon_{\mathrm{mach}}\approx2.2\times10^{-16}$.

\subsection{The stability bound}
\label{sec:cfl}

The bound follows from the dispersion relation of the scheme, which we need again in Section~\ref{sec:dispersion}
and therefore state once. Substituting a Fourier mode $\mathbf{e}^{n}_i\propto\exp(\mathrm{i}\mathbf{k}\cdot
\mathbf{x}_i-\mathrm{i}\omega n\Delta t)$ into \eqref{eq:leapfrog-h}--\eqref{eq:leapfrog-e} on a uniform grid of
pitch $h$, and eliminating $\mathbf{h}$, gives
\begin{equation}
  \omega^2 = \frac{4c_0^2}{h^2}\sum_{a\in\{x,y,z\}}\sin^2\Bigl(\frac{k_a h}{2}\Bigr),
  \label{eq:dispersion}
\end{equation}
the classical relation of the Yee lattice \cite{yee1966,taflove2005}. The factor four is what makes the
long-wavelength limit $\omega\to c_0|\mathbf{k}|$; with it, \eqref{eq:dispersion} reduces to the continuum
relation as $h\to0$ and the scheme is consistent.

Stability of the leapfrog requires $\omega_{\max}\Delta t\le2$, and \eqref{eq:dispersion} is largest at the corner
of the Brillouin zone, $k_a h=\pi$, where $\omega_{\max}=2c_0\sqrt{3}/h$. Hence
\begin{equation}
  \Delta t \le \sigma\,\frac{h_{\min}}{c_0\sqrt{3}}, \qquad \sigma=0.99 \text{ by default},
  \label{eq:cfl}
\end{equation}
the three-dimensional cubic-lattice Courant bound, with a safety factor we keep strictly below one.

We checked that \eqref{eq:cfl} is sharp rather than trusting it, because a bound that is merely sufficient is a
bound one pays for on every run. Table~\ref{tab:cfl} reports the outcome at three multiples of the bound.

\begin{table}[H]
  \centering
  \caption{Behaviour of the stepper at three multiples of the Courant bound \eqref{eq:cfl}, on a uniform grid
  driven from a random initial field. The bound is not conservative: a two per cent excess diverges.}
  \label{tab:cfl}
  \begin{tabular}{@{}lll@{}}
    \toprule
    $\Delta t/\Delta t_{\mathrm{CFL}}$ & outcome & remark \\
    \midrule
    $0.98$ & stable    & field bounded over the whole run \\
    $1.01$ & stable    & field bounded, as the bound predicts \\
    $1.02$ & divergent & growth to $6.7\times10^{122}$ over 5000 steps \\
    \bottomrule
  \end{tabular}
\end{table}

That the failure at $1.02$ is a clean overflow rather than a slow wobble is worth noting: the unstable mode grows
geometrically, so 5000 steps amplify it by the factor quoted. A code that merely looked at the first hundred steps
would have called it stable.

The structural point is that \eqref{eq:cfl} involves $h_{\min}$ and not an average. Every mode is bounded by the
worst cell, so a single small cell anywhere in the domain limits the step for the whole grid, and a mesh refined
locally to resolve a narrow gap pays for that refinement everywhere. This is why the examples keep the grid
uniform, why grading is attractive only when it is mild, and why local time stepping belongs on the list of things
to add rather than being a refinement of the present scheme.

\subsection{Numerical dispersion}
\label{sec:dispersion}

Equation \eqref{eq:dispersion} is exact for the scheme, so it also predicts how fast a wavefront travels. Writing
$v_{\mathrm{p}}=\omega/|\mathbf{k}|$ and using \eqref{eq:dispersion},
\begin{equation}
  \frac{v_{\mathrm{p}}}{c_0} = \frac{1}{|\mathbf{k}|h}
  \left[4\sum_{a}\sin^2\Bigl(\frac{k_a h}{2}\Bigr)\right]^{1/2},
  \label{eq:phase}
\end{equation}
which is strictly less than one for every non-zero $\mathbf{k}$ and tends to one as $|\mathbf{k}|h\to0$. On the
$40\times6\times6$ grid of the demonstration, with $h=0.02$ m cells and a pulse whose nominal wavelength spans
about ten cells, the measured wavefront speed is $0.904\,c_0$.

Equation \eqref{eq:phase} is worth evaluating at that resolution, because it shows both the sign and the size of
the effect. For propagation along an axis at ten cells per wavelength it gives $0.984\,c_0$; the measured deficit
is larger than that, and corresponds instead to about four cells per wavelength, where the same formula gives
$0.900\,c_0$. That is the resolution at which the leading edge of a broadband Ricker pulse actually sits, well
below the nominal wavelength of its spectrum. We therefore report the measurement as a dispersion measurement of
the tracked front rather than of the nominal wavelength, and note that the relation accounts for its size and,
more importantly, for its sign.

What matters for the credibility of the whole scheme is that the deficit is a smooth, direction-dependent,
second-order bias: it follows \eqref{eq:phase} to within the ambiguity of what the front is, it vanishes as
$h/\lambda\to0$, and it is not an instability and not a symptom of a wrong operator. It is the standard price of
the staggered grid \cite{taflove2005}, and it is one of the reasons the validation section compares frequencies
rather than propagation speeds.

\subsection{Where the two limits come from}

Both limits in this section are consequences of the one approximate step identified in
Section~\ref{sec:theory}, the replacement of the constitutive law by the diagonal star \eqref{eq:stars}. The
dispersion \eqref{eq:phase} is what a lumped material operator does to the phase of a short wave, and the bound
\eqref{eq:cfl} is the price of the explicit step that the same lumping buys. Neither can be removed by a better
integrator: the identity \eqref{eq:invariant} is already the best a fixed step can do, and it holds to rounding.
The next section turns to the frequency domain, where the same two operators give a generalised eigenproblem, and
where the structure that makes the time-domain path exact leaves behind a null space that has to be dealt with
before a spectrum can be believed.
